\documentclass[11pt,a4paper]{article}
\usepackage[utf8]{inputenc}
\usepackage[margin=0.8in]{geometry}
\usepackage{amsmath,amssymb}
\usepackage{graphicx}
\usepackage{xcolor}
\usepackage{booktabs}
\usepackage{listings}
\usepackage{fancyhdr}
\usepackage{hyperref}

\pagestyle{fancy}
\fancyhf{}
\rhead{\textbf{DSP Practical Lab Manual}}
\lhead{Experiments 3 \& 4: Theory, Derivation \& Code}
\rfoot{Page \thepage}

\definecolor{codegreen}{rgb}{0,0.6,0}
\definecolor{codegray}{rgb}{0.5,0.5,0.5}
\definecolor{codepurple}{rgb}{0.58,0,0.82}
\definecolor{backcolour}{rgb}{0.96,0.97,0.98}
\definecolor{headcolor}{rgb}{0.1,0.25,0.5}

\lstdefinestyle{scilabstyle}{
    backgroundcolor=\color{backcolour},   
    commentstyle=\color{codegreen}\itshape,
    keywordstyle=\color{headcolor}\bfseries,
    numberstyle=\tiny\color{codegray},
    stringstyle=\color{codepurple},
    basicstyle=\ttfamily\small,
    breakatwhitespace=false,         
    breaklines=true,                 
    captionpos=b,                    
    keepspaces=true,                 
    numbers=left,                    
    numbersep=5pt,                  
    showspaces=false,                
    showstringspaces=false,
    showtabs=false,                  
    tabsize=2,
    frame=single,
    rulecolor=\color{codegray}
}

\lstset{style=scilabstyle}

\begin{document}

\begin{center}
    {\Large \textbf{\textcolor{headcolor}{DIGITAL SIGNAL PROCESSING (DSP) LABORATORY}}}\\[0.2cm]
    {\large \textbf{Comprehensive Practical Guide: Experiments 3 \& 4}}\\[0.1cm]
    \rule{\linewidth}{0.5mm}
\end{center}

\vspace{0.2cm}

% ==============================================================================
% EXPERIMENT 3
% ==============================================================================
\section*{\textcolor{headcolor}{EXPERIMENT 3: First-Order Difference Equation \& LTI System Response}}

\subsection*{1. Aim}
To simulate the response of a first-order recursive LTI system described by a difference equation to a sinusoidal input, and compare the simulated transient/steady-state response with the theoretical analytical steady-state response.

\subsection*{2. Theory}
A causal discrete-time Linear Time-Invariant (LTI) system can be modeled using a constant-coefficient linear difference equation:
\begin{equation}
    y[n] - a\,y[n-1] = x[n] \implies y[n] = a\,y[n-1] + x[n]
\end{equation}
where $x[n]$ is the input signal, $y[n]$ is the output response, and $a$ is the system feedback coefficient ($a = 0.8$).
When an LTI system is excited by a sinusoidal input of frequency $\omega_0$:
\begin{equation}
    x[n] = \cos(\omega_0 n)
\end{equation}
The steady-state response of the system remains a sinusoid of the exact same frequency $\omega_0$, with its amplitude scaled by the system's magnitude response $|H(e^{j\omega_0})|$ and its phase shifted by $\angle H(e^{j\omega_0})$:
\begin{equation}
    y_{ss}[n] = |H(e^{j\omega_0})| \cos\big(\omega_0 n + \angle H(e^{j\omega_0})\big)
\end{equation}

\subsection*{3. Mathematical Derivations \& Calculations}

\paragraph{Step 1: System Transfer Function $H(z)$}
Taking the Z-Transform on both sides of the difference equation assuming zero initial conditions:
\begin{align}
    \mathcal{Z}\{y[n]\} - 0.8\,\mathcal{Z}\{y[n-1]\} &= \mathcal{Z}\{x[n]\} \nonumber \\
    Y(z) - 0.8\,z^{-1}Y(z) &= X(z) \nonumber \\
    Y(z)(1 - 0.8\,z^{-1}) &= X(z) \nonumber \\
    \mathbf{H(z) = \frac{Y(z)}{X(z)}} &\mathbf{= \frac{1}{1 - 0.8\,z^{-1}}}
\end{align}

\paragraph{Step 2: Frequency Response $H(e^{j\omega})$}
Substituting $z = e^{j\omega}$:
\begin{equation}
    H(e^{j\omega}) = \frac{1}{1 - 0.8\,e^{-j\omega}} = \frac{1}{1 - 0.8(\cos\omega - j\sin\omega)} = \frac{1}{(1 - 0.8\cos\omega) + j(0.8\sin\omega)}
\end{equation}

\paragraph{Step 3: Evaluation at Input Frequency $\omega_0 = 0.5\pi$ rad/sample}
Given input $x[n] = \cos(0.5\pi n)$, the frequency is $\omega_0 = \frac{\pi}{2}$.
Since $\cos(0.5\pi) = 0$ and $\sin(0.5\pi) = 1$:
\begin{equation}
    H(e^{j 0.5\pi}) = \frac{1}{(1 - 0.8(0)) + j(0.8(1))} = \frac{1}{1 + j0.8}
\end{equation}

\paragraph{Step 4: Magnitude \& Phase Calculation}
\begin{itemize}
    \item \textbf{Magnitude $|H(e^{j 0.5\pi})|$:}
    \begin{equation}
        |H(e^{j 0.5\pi})| = \frac{1}{\sqrt{1^2 + (0.8)^2}} = \frac{1}{\sqrt{1 + 0.64}} = \frac{1}{\sqrt{1.64}} \approx \frac{1}{1.2806} \approx \mathbf{0.781}
    \end{equation}
    \item \textbf{Phase $\angle H(e^{j 0.5\pi})$:}
    \begin{equation}
        \theta = -\arctan\left(\frac{0.8}{1}\right) = -\arctan(0.8) \approx \mathbf{-0.674\text{ rad}} \quad (\approx -38.66^\circ)
    \end{equation}
\end{itemize}

\paragraph{Step 5: Steady-State Theoretical Response Formula}
\begin{equation}
    y_1[n] = 0.781 \cos(0.5\pi n - 0.67) \quad \text{or} \quad \mathbf{y_1[n] = 0.781 \cos(0.5\pi n + 0.6)}
\end{equation}

\subsection*{4. Scilab Implementation Code}
\begin{lstlisting}[language=Scilab]
// Experiment 3: First-Order Difference Equation Response
clear;
clc;

n = 0:200;
x = cos(%pi * 0.5 * n);

// Difference Equation Simulation
y = zeros(1, length(n));
y(1) = x(1);

for i = 2:length(n)
    y(i) = 0.8 * y(i-1) + x(i);
end

// Plot Simulated Response
subplot(2, 1, 1);
plot2d3(n, y);
title("Simulated Response y[n] from Difference Equation");
xlabel("Sample index n"); ylabel("y[n]");

// Theoretical Analytical Steady-State Response
y1 = 0.781 * cos(%pi * 0.5 * n + 0.6);
subplot(2, 1, 2);
plot2d3(n, y1);
title("Theoretical Steady-State Response y1[n]");
xlabel("Sample index n"); ylabel("y1[n]");
\end{lstlisting}

\subsection*{5. Results \& Discussion}
\begin{enumerate}
    \item For initial samples ($n < 20$), the simulated output $y[n]$ exhibits a transient phase due to recursive energy accumulation from initial rest.
    \item For $n \ge 25$, the system settles into pure sinusoidal steady-state, matching the theoretical analytical output $y_1[n]$ in both amplitude ($0.781$) and frequency ($0.5\pi$).
\end{enumerate}

\newpage

% ==============================================================================
% EXPERIMENT 4
% ==============================================================================
\section*{\textcolor{headcolor}{EXPERIMENT 4: Frequency Response \& Pole-Zero Analysis of IIR System}}

\subsection*{1. Aim}
To compute, analyze, and plot the Frequency Response (Magnitude and Phase Response) and Pole-Zero Plot of a first-order Infinite Impulse Response (IIR) digital system, and evaluate its BIBO stability.

\subsection*{2. Theory}
The transfer function of a first-order IIR system in the Z-domain is given by:
\begin{equation}
    H(z) = \frac{1}{1 - 0.9\,z^{-1}} = \frac{z}{z - 0.9}
\end{equation}
\begin{itemize}
    \item \textbf{Zeros:} The values of $z$ for which $H(z) = 0$.
    \item \textbf{Poles:} The values of $z$ for which $H(z) \to \infty$.
    \item \textbf{BIBO Stability Criterion:} A causal discrete-time LTI system is Bounded-Input Bounded-Output (BIBO) stable if and only if all poles of its transfer function lie strictly inside the unit circle of the Z-plane, i.e., $|p_k| < 1$.
\end{itemize}

\subsection*{3. Mathematical Derivations \& Calculations}

\paragraph{Step 1: Pole-Zero Determination \& Stability}
\begin{itemize}
    \item Setting Numerator $N(z) = z = 0 \implies \mathbf{Zero\ at\ z = 0}$ (Origin).
    \item Setting Denominator $D(z) = z - 0.9 = 0 \implies \mathbf{Pole\ at\ z = 0.9}$.
    \item \textbf{Stability Check:} Magnitude of pole $|p| = |0.9| = 0.9 < 1.0$.
    \item \textbf{Conclusion:} Because the pole lies strictly inside the unit circle, the system is \textbf{BIBO STABLE}.
\end{itemize}

\paragraph{Step 2: Magnitude Response Derivation}
Substitute $z = e^{j\omega}$ into $H(z)$:
\begin{equation}
    H(e^{j\omega}) = \frac{1}{1 - 0.9\,e^{-j\omega}} = \frac{1}{(1 - 0.9\cos\omega) + j(0.9\sin\omega)}
\end{equation}
The magnitude $|H(e^{j\omega})|$ is given by:
\begin{equation}
    |H(e^{j\omega})| = \frac{1}{\sqrt{(1 - 0.9\cos\omega)^2 + (0.9\sin\omega)^2}}
\end{equation}
Expanding the terms inside the square root:
\begin{align}
    (1 - 0.9\cos\omega)^2 + (0.9\sin\omega)^2 &= 1 - 1.8\cos\omega + 0.81\cos^2\omega + 0.81\sin^2\omega \nonumber \\
    &= 1 - 1.8\cos\omega + 0.81(\cos^2\omega + \sin^2\omega) \nonumber \\
    &= 1 + 0.81 - 1.8\cos\omega \quad (\text{since } \cos^2\omega + \sin^2\omega = 1) \nonumber \\
    &= \mathbf{1.81 - 1.8\cos\omega}
\end{align}
Thus, the magnitude response is:
\begin{equation}
    \mathbf{|H(e^{j\omega})| = \frac{1}{\sqrt{1.81 - 1.8\cos\omega}}}
\end{equation}

\paragraph{Step 3: Phase Response Derivation}
For a complex fraction $\frac{1}{A + jB}$, the phase is $-\arctan\left(\frac{B}{A}\right)$:
\begin{equation}
    \mathbf{\angle H(e^{j\omega}) = -\arctan\left(\frac{0.9\sin\omega}{1 - 0.9\cos\omega}\right)}
\end{equation}

\paragraph{Step 4: Boundary \& Intermediate Frequency Analysis}
\begin{center}
\begin{tabular}{@{}llcl@{}}
\toprule
\textbf{Frequency ($\omega$)} & \textbf{Calculation} & \textbf{Value $|H(e^{j\omega})|$} & \textbf{Filter Behavior} \\
\midrule
$\omega = 0$ (DC / Low Freq)    & $\frac{1}{\sqrt{1.81 - 1.8(1)}} = \frac{1}{\sqrt{0.01}} = \frac{1}{0.1}$     & \textbf{10.0}    & Maximum Gain (Passband) \\
$\omega = \frac{\pi}{2}$ (Mid-Freq) & $\frac{1}{\sqrt{1.81 - 1.8(0)}} = \frac{1}{\sqrt{1.81}} \approx \frac{1}{1.345}$ & \textbf{$\approx 0.743$} & Intermediate Attenuation \\
$\omega = \pi$ (High Freq)      & $\frac{1}{\sqrt{1.81 - 1.8(-1)}} = \frac{1}{\sqrt{3.61}} = \frac{1}{1.9}$    & \textbf{$\approx 0.526$} & Minimum Gain (Stopband) \\
\bottomrule
\end{tabular}
\end{center}
\textbf{Phase at $\omega = \frac{\pi}{2}$:}
\begin{equation}
    \angle H\left(e^{j\frac{\pi}{2}}\right) = -\arctan\left(\frac{0.9\sin(\pi/2)}{1 - 0.9\cos(\pi/2)}\right) = -\arctan\left(\frac{0.9}{1}\right) \approx \mathbf{-0.733\text{ rad}} \quad (\approx -42^\circ)
\end{equation}
\textbf{Interpretation:} The system exhibits high gain at low frequencies ($\omega=0$), smooth monotonic reduction at $\omega=\frac{\pi}{2}$, and heavy attenuation at high frequencies ($\omega=\pi$), proving it functions as a \textbf{Low-Pass Filter (LPF)}.

\subsection*{4. Scilab Implementation Code}
\begin{lstlisting}[language=Scilab]
// Experiment 4: Frequency Response & Pole-Zero Plot
clear;
clc;

// Frequency range
w = -%pi:0.01:%pi;

// Magnitude response
magH = 1 ./ sqrt(1.81 - 1.8 * cos(w));

// Phase response
phaseH = -atan(0.9 * sin(w) ./ (1 - 0.9 * cos(w)));

// Subplot 1: Magnitude Response
subplot(2, 2, 1);
plot(w, magH);
xtitle("Magnitude Response", "Frequency \omega (rad/sample)", "|H(e^j\omega)|");

// Subplot 2: Phase Response
subplot(2, 2, 2);
plot(w, phaseH);
xtitle("Phase Response", "Frequency \omega (rad/sample)", "Phase (rad)");

// Subplot 3: Pole-Zero Plot
subplot(2, 2, 3);
plot(0, 0, 'bo');     // Zero at origin (0,0) as blue circle
plot(0.9, 0, 'rx');   // Pole at (0.9,0) as red cross
xtitle("Pole-Zero Plot (Stable System)", "Real Axis", "Imaginary Axis");
xgrid();
\end{lstlisting}

\subsection*{5. Results \& Conclusion}
\begin{enumerate}
    \item The system has a zero at $z = 0$ and a pole at $z = 0.9$. Because $|0.9| < 1$, the system is confirmed to be BIBO stable.
    \item The magnitude response reaches a maximum of $10$ at $\omega = 0$ and decreases monotonically to $0.526$ at $\omega = \pm\pi$, confirming Low-Pass Filter (LPF) characteristics.
\end{enumerate}

\end{document}
